% Peekaboo-VL report: 02_related. Included from ../peekaboo_report.tex.
\section{Related Work}
\label{sec:related}

\textbf{Video language models.} Qwen2.5-VL~\cite{qwen25vl2025} tells the
model the time of each frame through the way it numbers token positions
(a ``rotary position'' scheme). Qwen3-VL~\cite{qwen3vl2025} replaces this
with plain text timestamps, such as ``\texttt{<3.0 seconds>}'', placed
before each pair of frames. Qwen3.5~\cite{qwen35} reuses the Qwen3-VL
vision part on a new language model. In that model, three out of every four
layers are ``Gated DeltaNet''~\cite{yang2025gdn} layers, which work like a
running memory and are cheaper than standard attention; every fourth layer
is standard attention. We build on Qwen3.5 and use its text timestamps as
one of the things we ``tell'' the model.

\textbf{Question answering on drone and remote-sensing video.}
ReVA~\cite{yao2026reva} introduces the benchmark and the ReMoSense model.
ReMoSense learns camera-motion tokens from cost volumes (a dense map of how
well each patch matches the next frame) and object tokens from a recurrent
transformer. Our motion hint carries the same kind of information, but as
text, with no new model parts. EgoSIS~\cite{yang2026egosis} is the closest
drone work we know of. It fits a robust image-plane transform between frames
from optical flow, which is almost the same measurement our probe makes, and
turns it into four learned ``spatial evidence'' tokens per slice of video on
top of Qwen3-VL. On SIS-Bench this lifts an 8B model from 63.1\% to 76.2\%.
So the two routes start from the same signal and differ only in how it
reaches the language model: as learned tokens, or as a sentence.
UrbanVideo-Bench~\cite{zhao2025urbanvideo}
is another drone-video benchmark. Its videos are filmed from a drone flying
through city streets, and its questions test whether a model can follow
the flight, remember what it saw, and plan. Its 5,355 multiple-choice
questions have between two and seven options. Unlike ReVA, it asks mainly
about the drone's own movement and route, which is close to our
camera-motion questions. We use it only as an outside test of whether our
motion hint carries over to other drone videos; we never train on it.

\textbf{Camera motion in video models.} CameraBench~\cite{lin2025camerabench}
shows that classic geometry methods (``structure from motion'') fail on
camera movements that need understanding of the scene, while video models
fail on simple geometric ones. Combining the two works better than either.
Feng et al.~\cite{feng2026geocam} look inside a Qwen2.5-VL vision encoder
and find that camera motion is only weakly represented, especially in its
later layers. They extract camera motion with a 3D geometry model, turn it
into motion labels with a small classifier, and give these labels to the
video model in the prompt. This improves camera-motion answers without
fine-tuning the video model. Their approach is the closest to ours: both
\emph{tell} the model the camera motion in text. Ours differs in three
ways. We use a much simpler and faster estimator, because drone scenes are
nearly flat. We write the hint in the exact words of the benchmark's
answers. And we fine-tune the model together with the hint, so it learns
when to trust it. Two recent papers learn the motion instead.
CameraVLM~\cite{cheng2026cameravlm} trains camera and optical-flow tokens
together with the language model and argues that the motion has to be read
by the language model, not estimated on its own. CamDistill~\cite{du2026camdistill}
distills the camera tokens of a 3D model into the video model during
training, so the 3D model is not needed at test time. It also tries our kind
of shortcut, writing the 3D model's per-frame pose as text in the prompt,
without fine-tuning. That helps the model find \emph{when} the motion
changes but not \emph{which} motion it is. Our results point at the wording
as one reason. Given to a small model without fine-tuning, the same motion
numbers dropped camera-motion accuracy from 68.2\% to 55.7\% when written in
image terms, and brought it to 71.6\% when written in the answers' words
(Sec.~\ref{sec:cue}). Our estimator itself uses ORB
features~\cite{rublee2011orb}, simple matching, and a RANSAC
fit~\cite{fischler1981ransac}.

\textbf{Biases in multiple-choice answers.} Language models prefer some
option letters regardless of the content, for example ``A'' over ``D''.
PriDe~\cite{zheng2024pride} measures this preference by showing the same
question with the options in rotated orders, and then divides it out,
without needing any correct answers. Ok et al.~\cite{ok2026optionbias}
split this kind of bias into two parts: how often an answer flips when the
options are reordered, and how strongly the answers lean toward some
letters. Prior correction such as PriDe mostly fixes the lean, while
averaging over option orders mostly reduces the flips. We use both: PriDe
and option shuffling at test time, and random option order during training.

A related, older problem is the \emph{language prior}: patterns in the
questions and answers that reveal the answer without the image. Goyal et
al.~\cite{goyal2017vqav2} showed how strong such priors were in image
question answering, and Chakraborty et al.~\cite{chakraborty2026visualoptions}
show that video models still lean on the wording of text options: one strong
model drops from about 74\% to about 30\% when the same options are shown as
frames instead. For aerial video in particular, a recent meta-analysis of 67
papers~\cite{angelidis2026aerialgap} found that none of the aerial
question-answering benchmarks reports a blind (text-only) score, so there is
no floor to read their accuracies against. Our trained blind baseline
(Sec.~\ref{sec:blind}) gives that floor for ReVA.
